\documentclass[a4paper]{article}
\usepackage[square,sort,comma,numbers]{natbib}
\usepackage{blindtext} % Package to generate dummy text
\usepackage{charter} % Use the Charter font
\usepackage[utf8]{inputenc} % Use UTF-8 encoding
\usepackage[T1]{fontenc}
\usepackage{mathpazo}
\usepackage{microtype} % Slightly tweak font spacing for aesthetics
\usepackage{amsthm, amsmath, amssymb, amsfonts} % Mathematical typesetting
\usepackage{float} % Improved interface for floating objects
\usepackage{hyperref} % For hyperlinks in the PDF
\usepackage{graphicx, multicol} % Enhanced support for graphics
\usepackage{xcolor} % Driver-independent color extensions
\usepackage{pseudocode} % Environment for specifying algorithms in a natural way
\usepackage{datetime} % Uses YEAR-MONTH-DAY format for dates
\usepackage{mathabx}

\usepackage{tikz}
\usetikzlibrary{backgrounds}
\usetikzlibrary{arrows}
\usetikzlibrary{shapes,shapes.geometric,shapes.misc}
\usetikzlibrary{positioning}

% this style is applied by default to any tikzpicture included via \tikzfig
\tikzstyle{tikzfig}=[baseline=-0.25em,scale=0.5,node distance=1.5cm,
  variablenode/.style={circle, draw=black, thick, minimum size=5mm},
  observedvariablenode/.style={circle, draw=black, fill=gray!50, thick, minimum size=5mm},
  factornode/.style={rectangle, draw=black, thick, fill=black, minimum size=4mm},
  connection/.style={-, >=stealth, thick, rounded corners},
  message/.style={->, >=stealth, thick, rounded corners}]

% these are dummy properties used by TikZiT, but ignored by LaTex
\pgfkeys{/tikz/tikzit fill/.initial=0}
\pgfkeys{/tikz/tikzit draw/.initial=0}
\pgfkeys{/tikz/tikzit shape/.initial=0}
\pgfkeys{/tikz/tikzit category/.initial=0}

% standard layers used in .tikz files
\pgfdeclarelayer{edgelayer}
\pgfdeclarelayer{nodelayer}
\pgfsetlayers{background,edgelayer,nodelayer,main}

% style for blank nodes
\tikzstyle{none}=[inner sep=0mm]

% include a .tikz file
\newcommand{\tikzfig}[1]{%
{\tikzstyle{every picture}=[tikzfig]
\IfFileExists{#1.tikz}
  {\input{#1.tikz}}
  {%
    \IfFileExists{./figures/tikz/#1.tikz}
      {\input{./figures/#1.tikz}}
      {\tikz[baseline=-0.5em]{\node[draw=red,font=\color{red},fill=red!10!white] {\textit{#1}};}}%
  }}%
}

% the same as \tikzfig, but in a {center} environment
\newcommand{\ctikzfig}[1]{%
\begin{center}\rm
  \tikzfig{#1}
\end{center}}

% fix strange self-loops, which are PGF/TikZ default
\tikzstyle{every loop}=[]

\addtolength{\hoffset}{-2.25cm}
\addtolength{\textwidth}{4.5cm}
\addtolength{\voffset}{-3.25cm}
\addtolength{\textheight}{5cm}

\setlength{\parskip}{1ex}
\setlength{\parindent}{0in}

\DeclareUnicodeCharacter{03B1}{$\alpha$}

\usepackage{fancyhdr} % Headers and footers
\pagestyle{fancy} % All pages have headers and footers
\fancyhead{}\renewcommand{\headrulewidth}{0pt} % Blank out the default header
\fancyfoot[L]{} % Custom footer text
\fancyfoot[C]{} % Custom footer text
\fancyfoot[R]{\thepage} % Custom footer text
\newtheorem{thm}{Theorem}
\newtheorem{cor}{Corollary}

% define new commands
\newcommand{\Real}{{\mathbb R}}
\newcommand{\Normal}[3]{{\mathcal N} \left({#1};{#2},{#3}\right)}
\newcommand{\Gauss}[3]{{\mathcal G} \left({#1};{#2},{#3}\right)}
\newcommand{\NormalStandard}[1]{{\mathcal N} \left({#1}\right)}
\newcommand{\GaussStandard}[1]{{\mathcal G} \left({#1}\right)}
\newcommand{\NormalCDF}[3]{\Phi \left({#1};{#2},{#3}\right)}
\newcommand{\NormalStandardCDF}[1]{\Phi \left({#1}\right)}
\newcommand{\Bernoulli}[2]{{\mathrm{Ber}} \left({#1};{#2}\right)}
\newcommand{\Ber}[2]{{\mathcal B} \left({#1};{#2}\right)}
\newcommand{\expect}[1]{{\mathbb E \left[ {#1} \right]}}
\newcommand{\bs}[1]{{\boldsymbol{#1}}}
\newcommand{\var}[1]{{\mathbb V \left[ {#1} \right]}}
\newcommand{\dirac}[1]{{\delta \left( {#1} \right)}}
\newcommand{\diag}[1]{{\mathrm{diag}\left( {#1} \right)}}
\newcommand{\uniform}[1]{{\mathcal U} \left( {#1} \right)}
\newcommand{\identity}[1]{{{\mathbb{I}} \left( {#1} \right)}}
\newcommand{\incoming}[1]{\widecheck{#1}}
\newcommand{\outgoing}[1]{\widehat{#1}}
\newcommand{\intd}[1]{\ \mathrm{d}{#1}}
\newcommand{\transpose}[1]{{#1}^\top}
\newcommand{\kl}[3]{{\mathrm{KL}_{#1}} \left[{#2};{#3}\right]}
% \newcommand{\incoming}[1]{\stackrel{\rightarrow}{#1}}
% \newcommand{\outgoing}[1]{\stackrel{\leftarrow}{#1}}
\newcommand{\neighbours}[1]{{\mathrm{ne} \left( {#1} \right)}}


% define new environments
\newtheorem{theorem}{Theorem}
\newtheorem{lemma}{Lemma}

\theoremstyle{definition}
\newtheorem{definition}{Definition}

%----------------------------------------------------------------------------------------

\begin{document}

%-------------------------------
%	TITLE SECTION (do not modify unless you really need to)
%-------------------------------
\fancyhead[C]{}
\hrule \medskip
\begin{minipage}{0.295\textwidth}
    \raggedright
    \hfill\\
    % \footnotesize
    % Start: 24.4.2023\\
    % Return: 19.5.2023 \hfill\\
    % \href{https://classroom.github.com/classrooms/124387260-hpi-artificial-intelligence-teaching-pml-classroom}{https://classroom.github.com}
\end{minipage}
\begin{minipage}{0.4\textwidth}
    \centering
    \large
    Information Theory\\
\end{minipage}
\begin{minipage}{0.295\textwidth}
    \raggedleft
    \hfill\\
\end{minipage}
\medskip\hrule
\bigskip
 
In this short note, we will derive some of the results of unit 12 on {\em information theory}. 
\section*{Background}
The distribution that we will most commonly use is the one-dimensional Gaussian distribution.
\begin{definition}[One-Dimensional Gaussian Distribution]
    Given parameters $\mu \in \Real$, $\sigma \in \Real^+$, $\tau \in \Real$, and $\rho \in \Real^+$, a one-dimensional random variable $X \in \Real$ is distributed according to a one-dimensional Gaussian distribution if the density has the form $\Normal{\cdot}{\mu}{\sigma^2}$ or $\Gauss{\cdot}{\tau}{\rho}$ with
    \begin{align}
        \Normal{x}{\mu}{\sigma^2} & := \frac{1}{\sqrt{2\pi}\cdot \sigma} \cdot \exp \left( -\frac{(x-\mu)^2}{2\sigma^2}\right) \,, \label{eq:Normal_definition} \\
        \Gauss{x}{\tau}{\rho}     & := \sqrt{\frac{\rho}{2\pi}} \cdot \exp\left(-\frac{\tau^2}{2\rho} \right) \cdot \exp \left(\tau\cdot x - \rho \cdot \frac{x^2}{2} \right) \,. \label{eq:Gauss_definition}
    \end{align}
    If $\mu = \tau = 0$ and $\sigma^2 = \rho = 1$ we simply write 
    \begin{align}
        \NormalStandard{x} & := \Normal{x}{0}{1} \,, \\
        \GaussStandard{x}  & := \Gauss{x}{0}{1} \,.
    \end{align}
\end{definition}
Note that both definitions are identical when using the identities
\begin{align}
    \tau & = \frac{\mu}{\sigma^2} & \mbox{and} &  & \rho     & = \sigma^{-2}\,,                                       &  &  & \mbox{or} \\
    \mu  & = \frac{\tau}{\rho}    & \mbox{and} &  & \sigma^2 & = \rho^{-1} \,. \label{eq:Gauss_Normal_transformation}
\end{align}
Also note that if $X \sim \Normal{\cdot}{\mu}{\sigma^2}$ we know that $\expect{X}=\mu$, $\expect{X^2}=\mu^2 + \sigma^2$, and $\var{X}=\sigma^2$.

\section*{Entropy}
\begin{definition}[Entropy]
    Given $b > 1$, the entropy $H_b[X]$ of a random variable $X$ with density $p_X$ is defined as
    \begin{align}
        H_b[X] & := -\expect{\log_b \left( p_X(X) \right)} \,.
    \end{align}
    If $X$ is a discrete random variable, we can write this as
    \begin{align}
        H_b(X) & = -\sum_{x} p_X(x) \cdot \log_b\left( p_X(x) \right)  \,.
    \end{align}
\end{definition}

\begin{theorem}[Maximumm Entropy]
    The entropy $H_b[X]$ of a discrete random variable $X$ is maximized if and only if $X$ is uniformly distributed, that is, if $p_X(x) = \frac{1}{|\mathcal{X}|}$ for all $x \in \mathcal{X}$.
\end{theorem}
\begin{proof}
    The proof follows from Jensen's inequality, which states that for any concave function $f$ and random variable $X$ we have $\expect{f(X)} \leq f(\expect{X})$. The function $f(x) = \log_b(x)$ is concave for $x > 0$, and thus we have
    \begin{align}
        H_b[X] & = -\expect{\log_b(p_X(X))} = \expect{\log_b\left(\frac{1}{p_X(X)}\right)} \leq \log_b\left(\expect{\frac{1}{p_X(X)}}\right) = \log_b(|\mathcal{X}|) \,.
    \end{align}
    Since $\log_b(\cdot)$ is a non-linear function, the equality holds if and only if $p_X(x)$ is constant, which is the case if and only if $X$ is uniformly distributed.
\end{proof}


\begin{theorem}[Entropy of a Gaussian Distribution]
    The natural entropy $H_e[X]$ of a random variable distributed according to a one-dimensional Gaussian distribution $\Normal{\cdot}{\mu}{\sigma^2}$ is given by
    \begin{align}
        H_e[X] & = \frac{1}{2}\left(\log_e(2\pi\sigma^2) + 1 \right) \,.
    \end{align}
\end{theorem}
\begin{proof}
    The proof follows from the definition of the entropy and the definition of the Gaussian distribution.
    \begin{align*}
        H_e[X] & = -\expect{\log_e\left(\Normal{X}{\mu}{\sigma^2}\right)} \\
               & = -\expect{\log_e\left(\frac{1}{\sqrt{2\pi\sigma^2}} \cdot \exp \left( -\frac{(X-\mu)^2}{2\sigma^2}\right)\right)} \\
               & = -\expect{-\frac{1}{2}\log_e(2\pi\sigma^2) - \frac{(X-\mu)^2}{2\sigma^2}}  \\ 
               & = -\left(-\frac{1}{2}\log_e(2\pi\sigma^2) - \frac{\sigma^2}{2\sigma^2}\right) = \frac{1}{2}\left(\log_e(2\pi\sigma^2) + 1 \right) \,.
    \end{align*}
    The last equality follows from the fact that $\expect{(X - \mu)^2} = \var{X} = \sigma^2$.
\end{proof}

\section*{KL Divergence}

\begin{definition}[KL Divergence]
    Given $b > 1$ and two random variables $X$ and $Y$ with densities $p_X$ and $p_Y$, the KL divergence $\kl{b}{X}{Y}$ is defined as
    \begin{align}
        \kl{b}{X}{Y} & := \int p_X(x) \cdot \log_b \left(\frac{p_X(x)}{p_Y(x)}\right) \intd{x} \,.
    \end{align}
\end{definition}

\begin{theorem}[KL Divergence of a Gaussian Distribution]
    The KL divergence $\kl{e}{X}{Y}$ of two random variables $X$ and $Y$ distributed according to a one-dimensional Gaussian distributions $\Normal{\cdot}{\mu_X}{\sigma_X^2}$ and $\Normal{\cdot}{\mu_Y}{\sigma_Y^2}$, respectively, is given by
    \begin{align}
        \kl{e}{X}{Y} & = \frac{1}{2}\left(\log_e\left(\frac{\sigma_Y^2}{\sigma_X^2}\right) + \frac{\sigma_X^2 + (\mu_X - \mu_Y)^2}{\sigma_Y^2} - 1\right) \,.
    \end{align}
\end{theorem}

\begin{proof}
    In the proof we use the following variance-decomposition identity for $X \sim \Normal{\cdot}{\mu_X}{\sigma_X^2}$:
    \begin{align}
        \expect{(X - c)^2} & =  \expect{((X - \mu_X) + (\mu_X - c))^2} = \expect{(X - \mu_X)^2} + \expect{(\mu_X - c)^2} = \sigma_X^2 + (\mu_X - c)^2\,, \label{eq:variance_decomposition}
    \end{align}
    where the third step follows from the fact that $\expect{(X - \mu_X)} = 0$. The proof follows from the definition of the KL divergence and the definition \eqref{eq:Normal_definition} for $X \sim \Normal{\cdot}{\mu_X}{\sigma_X^2}$.
    \begin{align*}
        \kl{e}{X}{Y} & = \expect{\log_e\left(\frac{p_X(X)}{p_Y(X)}\right)} \\
                     & = \expect{\log_e\left(\frac{\Normal{X}{\mu_X}{\sigma_X^2}}{\Normal{X}{\mu_Y}{\sigma_Y^2}}\right)} \\
                     & = \expect{\log_e\left(\frac{\sqrt{2\pi\cdot \sigma_Y^2}}{\sqrt{2\pi\cdot \sigma_X^2}} \cdot \exp\left(-\frac{(X-\mu_X)^2}{2\sigma_X^2} + \frac{(X-\mu_Y)^2}{2\sigma_Y^2}\right) \right)} \\
                     & = \expect{\frac{1}{2}\log_e\left(\frac{\sigma_Y^2}{\sigma_X^2}\right) - \frac{(X-\mu_X)^2}{2\sigma_X^2} + \frac{(X-\mu_Y)^2}{2\sigma_Y^2}} \\
                     & = \frac{1}{2} \cdot \left(\log_e\left(\frac{\sigma_Y^2}{\sigma_X^2}\right) - 1 + \frac{\sigma_X^2 + \left( \mu_X - \mu_Y \right)^2}{\sigma_Y^2}\right) \,,
    \end{align*}
    where we used the variance-decomposition identity \eqref{eq:variance_decomposition} in the penultimate step.
\end{proof}

\end{document}
